\section{Related work}\label{sec:related}

\paragraph{Stabilizer simulation.} The Gottesman Knill theorem makes Clifford circuits efficiently simulable through
the Heisenberg picture of stabilizer groups~\cite{gottesman1998heisenberg}. Aaronson and Gottesman gave the tableau
algorithm with destabilizers~\cite{aaronson2004chp}, and Anders and Briegel a graph state representation that they found
faster than the tableau algorithm~\cite{anders2006graph}. Stim~\cite{gidney2021stim} combines a tableau simulation of one
reference shot with Pauli frame sampling of all other shots, vectorised over shots with SIMD instructions; it also
derives the detector error model of a circuit, which decoders such as PyMatching consume~\cite{higgott2022pymatching,
higgott2025sparse}. SymPhase accumulates possible Pauli faults as symbolic phases during a single pass over the circuit and draws
samples by substituting values into the symbolic measurement outcomes; its authors report shorter sampling times than
Stim on their benchmarks~\cite{fang2024symphase}. All of these run on CPUs.

\paragraph{Accelerated stabilizer simulation.} Osama, Thanos and Laarman reformulate tableau evolution and
measurement for GPUs in the QuaSARQ simulator; they report speedups of up to 105 times on circuits with up to 180,000
qubits, consistently shorter run times than Stim, and an advantage in many shot sampling on large
workloads~\cite{osama2026gpu}. quEStab extends the tableau to mixed Clifford and non Clifford circuits on several
GPUs~\cite{shin2026questab}. These GPU simulators are built around the tableau of wide circuits; the sampler of this
paper instead streams the Pauli frames of many shots of a small QEC circuit through a fixed pipeline, and its output is
the detector and observable bits that a decoding study consumes. We are not aware of an FPGA implementation of Pauli
frame sampling other than our own first kernel (Section~\ref{sec:first}); searches of Crossref and arXiv for FPGA
based stabilizer, Clifford or Pauli frame simulators found none.

\paragraph{FPGAs for quantum circuit simulation.} FPGA emulation of quantum circuits dates back two
decades~\cite{khalid2004fpga}. Statevector simulators on HBM equipped cards reach 30 qubits on one
card~\cite{tiwari2026fpga30}, and our four card simulator distributes a statevector over the cards of one host and
treats each card as a node~\cite{ali2026multifpga}. Trefoil and Qu-Trefoil distribute much larger states over FPGAs
attached to SATA storage~\cite{niwase2023trefoil,wei2025qutrefoil}. Statevector simulation is limited by memory and by
the exchange of amplitudes between devices; Pauli frame sampling is limited by neither, which is why it scales over
nodes without communication.

\paragraph{FPGAs in quantum error correction.} Most FPGA work in QEC concerns decoding, where latency matters:
lookup table decoders for small codes~\cite{das2022lilliput}, distributed union find decoders~\cite{
delfosse2021unionfind,liyanage2023fpga,liyanage2024distributed}, HeliosNet, which decodes lattice surgery operations
on five FPGAs connected as a network~\cite{liyanage2024multifpga}, the Collision Clustering decoder on FPGA and ASIC
hardware with megahertz decoding rates~\cite{barber2025decoder}, and real time correction for the Shor and Steane
codes~\cite{ali2026qec}. These systems consume syndromes; the sampler of this
paper produces them, for the offline studies that design codes and test decoders. The two are complementary, and a
sampler that runs on the same class of device as a decoder can feed it test data at hardware rates.

\paragraph{High level synthesis.} The kernel is written in C++ and compiled with Vitis HLS~\cite{amd_ug1399}; Cong et
al.\ review the methodology~\cite{cong2011hls}. The design choices below (dataflow stages connected by streams, loops
pipelined to one iteration per cycle, arrays partitioned into registers) are standard HLS constructs; the
contribution is how the sampling algorithm is arranged so that these constructs reach full rate.
